\chapter{Pruebas, validación y resultados}
\label{cap:pruebas}

\section{Plan de pruebas}
\label{sec:plan-pruebas}

Las pruebas se organizan en tres niveles:

\begin{enumerate}
  \item \textbf{Unitarias:} codecs TCP, parseo UUID/NID, utilidades de gateway IP (\texttt{PctFrameCodecTest}, \texttt{PctNidTest}, \texttt{PctNetworkHelperTest}).
  \item \textbf{Integración L2:} handshake HELLO/JOIN en dos dispositivos; verificación UUID en topología.
  \item \textbf{Sistema (EXP-01..06):} escenarios multisalto en laboratorio con ADB/logcat.
\end{enumerate}

\section{EXP-01: Interconexión dos GO}
\label{sec:resultados-exp01}

\textbf{Configuración:} D1 ROOT (Samsung A24), D2 hijo (Vivo Y21s).

\textbf{Procedimiento:} D1 \texttt{start()} $\rightarrow$ espera anuncio \texttt{\_pct-ctrl} $\rightarrow$ D2 \texttt{start()}.

\textbf{Criterios:}
\begin{itemize}
  \item D2 asocia STA al SSID de D1.
  \item D2 mantiene GO propio (H1).
  \item TCP HELLO upstream completado; UUID padre visible en UI.
  \item $T_{\mathrm{join}} < 15$~s.
\end{itemize}

% TODO: completar tabla de resultados medidos en defensa
\begin{table}[htbp]
  \centering
  \caption{Resultados EXP-01 (plantilla --- completar con mediciones)}
  \label{tab:exp01}
  \begin{tabular}{@{}lcc@{}}
    \toprule
    \textbf{Métrica} & \textbf{Objetivo} & \textbf{Observado} \\
    \midrule
    STA asociada            & Sí  & --- \\
    GO hijo activo post-STA & Sí  & --- \\
    HELLO L2 completado     & Sí  & --- \\
    $T_{\mathrm{join}}$ (ms) & $<15000$ & --- \\
    \bottomrule
  \end{tabular}
\end{table}

\section{EXP-03: Cadena tres nodos}
\label{sec:resultados-exp03}

Topología: A (ROOT) $\leftarrow$ B (BRIDGE) $\leftarrow$ C (LEAF). C envía DATA a A; B reenvía por \texttt{ForwardWorker}. Criterio: entrega en C con \texttt{hop\_limit} decrementado y ruta ACTIVE en tablas intermedias.

\section{EXP-04 y EXP-05: Unión tardía y fallo BRIDGE}
\label{sec:exp04-exp05}

EXP-04 verifica que un nodo se una a red operativa sin \texttt{removeGroup} en nodos existentes. EXP-05 simula caída de B: hijos de B detectan NODE\_DOWN; B re-converge vía \texttt{\_pct-seek}; \texttt{session\_epoch} de conversaciones activas se mantiene (H5).

\section{Pruebas unitarias}
\label{sec:pruebas-unitarias}

Ejecución:
\begin{verbatim}
cd pct/library && ./gradlew :lib-pct-core:test
\end{verbatim}

Las pruebas validan round-trip de frames PCT1, normalización de UUID y resolución de gateway IPv4 en \texttt{Network} STA.

\section{Análisis de logs}
\label{sec:analisis-logs}

Tag \texttt{PctMesh} registra: fases bootstrap, candidatos DNS-SD, eventos STA, conexiones TCP (\texttt{local=}/\texttt{remote=}), HELLO upstream/downstream, errores self-connect.

Indicadores de fallo diagnosticados:
\begin{itemize}
  \item \texttt{pct=0 candidatos=0} con peers visibles: hijo escaneó antes del anuncio ROOT.
  \item \texttt{local=192.168.49.1 remote=192.168.49.1}: bind a GO propio en lugar de STA.
  \item UUID truncado con ceros: instance DNS sin merge TXT.
\end{itemize}

\section{Discusión}
\label{sec:discusion-resultados}

Los resultados confirman viabilidad de H1/H2 en hardware clase~$\alpha$. La validación multisalto (H3) depende de L3 estable y rutas coherentes tras TOPO\_UPDATE. Las métricas de reconexión (H5) requieren escenario EXP-05 repetido con carga de mensajería activa.

\section{Síntesis del capítulo}
\label{sec:sintesis-pruebas}

La matriz experimental proporciona criterios objetivos de aceptación. Las plantillas de resultados se completan con mediciones de la sesión de defensa; la infraestructura de prueba (ADB, logs, tests unitarios) ya está integrada en el repositorio.
